\chapter{Software Artifacts}\label{chapter:artifacts}

\paragraph{Prerequisites.} The reproduction recipes below assume a Linux 5.6+ kernel with cgroup~v2~\cite{cgroupv2} mounted at \texttt{/sys/fs/cgroup/}, \texttt{liburing} 2.0+ headers, GCC 11+, CMake 3.20+, and a user account with \texttt{sudo} for cgroup creation and kernel-cache drop operations. System swap must be disabled (\texttt{sudo swapoff -a}) so cgroup \texttt{MemorySwapMax=0} is enforceable end-to-end. The model file must be a GPT-OSS-20B or GPT-OSS-120B GGUF in F16 quantization~\cite{openai2025gptoss} from the \texttt{ggml-org/gpt-oss-*-GGUF} repositories on Hugging Face. The \texttt{.bscexp} sidecar is produced by the \texttt{bsc-pack-experts} preprocessor at first run (\autoref{sec:art-llamacpp}). The test machine described in \autoref{sec:eval-hardware} is the reference hardware; the methodology transfers to any host with NVMe-attached storage.


\section{Forked llama.cpp}\label{sec:art-llamacpp}

Repository: \texttt{github.com/ErsjanKeri/llama.cpp}, default branch \texttt{master}, tag \texttt{bsc-thesis-v1} (commit \texttt{5aec1115}). The fork carries twenty-six commits on top of upstream llama.cpp \cite{ggerganov2023llamacpp} covering the tensor tracer, the io\_uring + O\_DIRECT expert loader, the fused MoE custom op, and the three pipeline modes (projection-group overlap, async-projection-overlap, async-experts). Line-number citations below refer to the tagged commit.

\begin{table}[htbp]
\caption{Runtime flags introduced by the fork.}
\label{tab:runtime-flags}
\centering
\small
\begin{tabular}{@{}lp{8.0cm}@{}}
\toprule
Flag & Purpose \\
\midrule
\texttt{--pin-compute-weights} & \texttt{mlock} attention, output, and norm weights (excludes experts and embedding) \\
\texttt{--eager-compute} & \texttt{memset} compute buffer at \texttt{graph\_reserve} so its 413~MiB is physically resident \\
\texttt{--uring-experts} & enable the application-managed expert loader (requires \texttt{.bscexp}) \\
\texttt{--uring-cache-slots N} & cache size in slots; 0 disables the cache \\
\texttt{--uring-cache-policy} & eviction policy: \texttt{lru}, \texttt{lfu}, or \texttt{lfu-aging} \\
\texttt{--uring-aging-mult N} & LFU-aging period multiplier (default 10, tuned optimum 3) \\
\texttt{--uring-projection-overlap} & projection-group overlap (down async, up+gate sync) \\
\texttt{--uring-async-projection-overlap} & async-projection-overlap (split-tag, first-ready dispatch) \\
\texttt{--uring-async-experts} & async-experts (per-(expert, projection) tags, eager dispatch) \\
\texttt{--trace-mode} & tensor trace filter: \texttt{off}, \texttt{experts}, or \texttt{all} \\
\bottomrule
\end{tabular}
\end{table}

The three pipeline-selection flags are mutually exclusive; the constraint is enforced in \texttt{common/arg.cpp}.

\paragraph{Key code paths added.}\hspace{0pt}
\begin{itemize}\setlength{\itemsep}{0pt}
  \item \texttt{src/llama-io-uring.\{h,cpp\}}, \texttt{src/llama-io-uring-buf.\{h,cpp\}}: application-managed expert loader and cache buffer.
  \item \texttt{src/llama-moe-pipeline.\{h,cpp\}}: fused MoE FFN custom op for the overlap variants.
  \item \texttt{src/llama-graph.cpp}: graph integration, the phase callbacks, the \texttt{make\_compact} view fix, and the \texttt{tensor->extra} wiring.
  \item \texttt{ggml/src/ggml-cpu/ggml-cpu.c:1630--1644}: patched per-expert pointer indirection in \texttt{mul\_mat\_id}.
  \item \texttt{ggml/include/tensor\_trace.h}, \texttt{ggml/src/tensor\_trace.c}: tensor tracer.
  \item \texttt{tools/bsc-pack-experts/}: offline \texttt{.bscexp} pack-file builder.
  \item \texttt{tools/gguf-dump/gguf-dump.cpp}: MXFP4-aware metadata exporter.
\end{itemize}


\section{Tensor-Tracing Toolchain}\label{sec:art-tracing}

Repository directory: \texttt{\textasciitilde/BSC/tensor-tracing/}. The toolchain consumes the binary trace emitted by \texttt{ggml/src/tensor\_trace.c} (1024-byte cache-aligned records, see \autoref{chapter:tracer}) plus the GGUF metadata CSV from \texttt{llama-gguf-dump} and the per-token Graphviz \texttt{.dot} files from llama.cpp, and produces JSON payloads for two visualization frontends.

\begin{table}[htbp]
\caption{Python parsers in \texttt{tools/} of the tensor-tracing toolchain.}
\label{tab:python-parsers}
\centering
\small
\begin{tabular}{@{}lp{8.0cm}@{}}
\toprule
File & Purpose \\
\midrule
\texttt{parse\_trace.py} & 1024-byte binary trace records from \texttt{/tmp/tensor\_trace.bin} into per-token JSON; mirrors the C \texttt{ggml\_op} enum at lines 53--220 \\
\texttt{parse\_csv.py} & GGUF metadata CSV into \texttt{memory-map.json}; emits per-expert rows for the 2,691-tensor GPT-OSS-20B map \\
\texttt{parse\_dot.py} & per-token Graphviz \texttt{.dot} graphs into JSON nodes/edges, with layer-ID extraction for both \texttt{blk.N.} and \texttt{name-N} suffix patterns \\
\texttt{parse\_buffer\_stats.py} & \texttt{/tmp/buffer\_stats.jsonl} into an alloc/dealloc timeline plus per-buffer summary (ANY / WEIGHTS / COMPUTE) \\
\texttt{preprocess\_all.py} & orchestrator that runs the four parsers above and lays out \texttt{webui/public/data/} \\
\texttt{simulate\_from\_dump.py} & policy simulator: replays the access sequence dumped via \texttt{BSC\_CACHE\_DUMP} against LRU, LFU, LFU-aging, ARC, W-TinyLFU, and Belady. LRU and LFU-aging match the deployed C runtime to within rounding (34.13\% simulated vs 33.80\% measured at 7~GiB, c=250). \\
\bottomrule
\end{tabular}
\end{table}

\paragraph{Visualization frontends.} Two were built because no single technology stack handles both interactive 1--2 token introspection and 100+ token bulk analysis well.

The WebUI lives at \texttt{webui/}. Stack: React 19, Zustand 5 for state, Three.js 0.182 with \texttt{@react-three/fiber}/\texttt{drei} for the 3D heatmap, Cytoscape.js 3.33~\cite{cytoscape_js} with \texttt{cytoscape-dagre} for graph layout, \texttt{react-window} for virtual scrolling, Vite 7, TailwindCSS 3.4, TypeScript strict mode. It has four primary views (\texttt{GraphView}, \texttt{TraceView}, \texttt{HeatmapView}, \texttt{ViewContainer}) and is well-suited to 1--2 tokens (around 1,700 trace entries); beyond ten tokens, React reconciliation cost dominates.

The DesktopUI lives at \texttt{desktopui/}. Stack: C++17, OpenGL, GLFW3, vendored ImGui and nlohmann/json. \texttt{TraceTableView} reaches 110 FPS at 170,000+ rows via virtual scrolling, and \texttt{launch\_all\_domains.sh} opens one window per semantic domain. This is the workhorse for 100+ token analyses.

Both frontends consume identical JSON schemas, so a trace produced for one is loadable by the other.


\section{Experiment Harness}\label{sec:art-harness}

Repository directory: \texttt{\textasciitilde/BSC/time-tracking/}. The canonical runner is \texttt{run\_cgroup\_experiments.py}, which executes each iteration of \texttt{llama-completion} inside a fresh \texttt{bscexp.slice} cgroup-v2 scope created by \texttt{systemd-run}. The slice name is deliberately hyphen-free, because \texttt{bsc-experiment.slice} would land in \texttt{/sys/fs/cgroup/bsc.slice/bsc-experiment.slice/}: systemd treats hyphens as cgroup-hierarchy separators, and the diagnostic stats reader expects flat paths under \texttt{/sys/fs/cgroup/bscexp.slice/}.

{\sloppy
A bash wrapper inside the scope reads \texttt{memory.peak}, \texttt{memory.current}, \texttt{memory.swap.peak}, and \texttt{memory.events} from the still-existing cgroup before \texttt{systemd} tears the scope down, emitting \texttt{BSC\_CGROUP\_*=...} lines for the parser to scrape. Each iteration is preceded by \texttt{systemctl stop bscexp.slice}, \texttt{drop\_caches}, and a 5-second settle. All monitoring runs outside the cgroup so it never competes for the budget.
\par}

\paragraph{Configuration files.} Per-experiment settings live in JSON files alongside the harness. The two that drive the canonical sweep are \texttt{settings\_thesis\_final.json} (GPT-OSS-20B) and \texttt{settings\_thesis\_final\_120b\_retry.json} (GPT-OSS-120B).

\paragraph{Result archive layout.} Output lands in \texttt{results/cgroup\_<timestamp>/}, with one CSV per configuration plus the captured stdout of every iteration as \texttt{.log} files. Because \texttt{systemd-run} executes the inner command as root for cgroup setup, the result directories are root-owned; subsequent reads from user processes need either \texttt{sudo} or a follow-up \texttt{chown}. The harness does not chown automatically because doing so would race with downstream parsers in some setups.


\section{Datasets and Result Archives}\label{sec:art-datasets}

The thesis cites four classes of primary artifact. All paths are absolute under \texttt{\textasciitilde/BSC/}.

\paragraph{Research journal.} \texttt{journal/INDEX.md} is the annotated index, listing every dated entry with its status (\textsc{authoritative}, \textsc{reference only}, \textsc{deprecated-numbers}). Daily entries live as \texttt{journal/2026-MM-DD.md}; multi-day work folds into a dated directory with \texttt{summary.md}.

\paragraph{Cited result directories.} \texttt{time-tracking/results/} holds CSV output for the canonical sweep. The authoritative directories are \texttt{cgroup\_20260501\_191023/} (GPT-OSS-20B, 96 runs, 32 configurations $\times$ 3 iterations including the LRU/LFU/LFU-aging policy sweep at the tight 7~GiB cell) and \texttt{cgroup\_20260502\_140144/} (GPT-OSS-120B retry, 72 runs, 24 configurations $\times$ 3 iterations across the four cgroup budgets, LRU policy throughout because cache exceeds the per-token working set at every budget). The driving settings are \texttt{settings\_thesis\_final.json} and \texttt{settings\_thesis\_final\_120b\_retry.json}. The corresponding deterministic generation outputs at the canonical seed=42 are stored as \texttt{generation\_outputs/\{20b,120b\}\_multitopic\_seed42.txt}.


\section{Reproducing the Headline Numbers}\label{sec:art-reproducing}

The canonical sweep at \texttt{cgroup\_20260501\_191023/} (GPT-OSS-20B) and \texttt{cgroup\_20260502\_140144/} (GPT-OSS-120B) produced the headline numbers cited in \autoref{chapter:evaluation}. Reproducing them takes five steps on the same hardware, with the same kernel and the same NVMe.

Step 1 clones the fork and checks out the canonical commit. The thesis state is the tag \texttt{bsc-thesis-v1} (commit \texttt{5aec1115}).

\begin{lstlisting}[basicstyle=\ttfamily\small,breaklines=true]
# 1. Clone the fork and check out the canonical tag.
git clone https://github.com/ErsjanKeri/llama.cpp ~/llama.cpp
cd ~/llama.cpp && git checkout bsc-thesis-v1

# 2. Build llama-completion and bsc-pack-experts.
cmake -B build -DCMAKE_BUILD_TYPE=Release
cmake --build build --target llama-completion bsc-pack-experts -j$(nproc)

# 3. Produce the .bscexp sidecar pack file for the model.
./build/bin/bsc-pack-experts <model.gguf> -o <model.gguf>.bscexp --verify

# 4. Switch to the harness directory.
cd ~/BSC/time-tracking

# 5a. Run the canonical 20B sweep.
sudo python3 run_cgroup_experiments.py \
    --settings settings_thesis_final.json

# 5b. Run the canonical 120B retry sweep.
sudo python3 run_cgroup_experiments.py \
    --settings settings_thesis_final_120b_retry.json
\end{lstlisting}

The harness drops the page cache between iterations, wraps each \texttt{llama-completion} invocation in a fresh \texttt{bscexp.slice} cgroup with \texttt{MemorySwapMax=0}, runs three iterations per configuration, and writes one CSV per configuration into \texttt{results/cgroup\_<timestamp>/}.

The CSVs should match the canonical numbers within the cited CV bounds: mean CV across the 168-run sweep is 0.23\%, with the noisiest cell at 2.29\% (20B 8~GiB async-experts). If they do not, the most likely cause is a checkout other than tag \texttt{bsc-thesis-v1}; verify the tag matches the appendix. The next most likely cause is page-cache contamination between iterations on a host where the harness was not run with \texttt{sudo}; \texttt{drop\_caches} requires root.
